\subsection{A native clock read at wave arrival}

The same80000-dimensional matter--geometry--wave source prepares the right
native cell in an equal coherent superposition of levels0 and1.
A retained setting bit selects either the $X$ or $Y$ clock instrument when
the wave arrives. All native complements, geometry outcomes and timeouts
are included. The resulting four active setting/sign frequencies determine
the clock coherence without assigning simultaneous sharp $X$ and $Y$ values.
Independently prepared vacuum-clock controls provide a strictly increasing
phase calibration for the sixteen event labels.

Using this measured event-to-cycle dictionary, the source-on minus
source-off mean arrival difference lies in
\[
 [0.0013506,\,0.0013547]\quad\text{native clock cycles}.
\]
The phase of the actual aggregate arrival-clock coherence gives a separate
positive response in
$[0.0004679,0.0004687]$ cycles. It is a circular phase statistic, rather than
the mean arrival delay. Complete finite numerical bounds control both
quantities. A finite-repetition theorem also certifies the sign of the
phase response with retained failure records; its large repetition count
is an analytic resource bound.

Within a terminal event, clock, field and geometry effects act on distinct
tensor factors. Their products commute, and every permitted ordering has
the same full joint record. Copying these orthogonal records supplies two
observer readouts without copying an arbitrary quantum state. This closes
an operational clock--wave comparison on one specified construction.
No physical-second calibration, Einstein proper-time identification or
fundamental global-time variable is assumed.
